\section{Dados e métodos}

\subsection{Conjuntos fornecidos}

Cada janela possui \(\valNWin\) amostras, adquiridas a \(\valFs\,\mathrm{Hz}\).
O template do pulso tem \(\valNTmpl\) amostras. Os arquivos utilizados são:

\begin{table}[H]
\centering
\caption{Arquivos de dados do Trabalho~2.}
\label{tab:dados-t2}
\begin{tabular}{llp{7.2cm}}
\toprule
Arquivo & Dimensão & Uso \\
\midrule
\texttt{template\_pulso.csv} & \(\valNTmpl\) amostras & Forma \(s[n]\) \\
\texttt{ruido\_treino.csv} & \(\valNNoise\) janelas & Estatística do ruído e limiar \\
\texttt{sinais\_treino.csv} & \(\valNTrain\) janelas & Ajuste e comparação de métodos \\
\texttt{rotulos\_treino.csv} & \(\valNTrain\) rótulos & \texttt{event\_present}, \texttt{t0\_sample}, \(A\) \\
\texttt{sinais\_teste.csv} & \(\valNTest\) janelas & Avaliação final \\
\texttt{rotulos\_teste.csv} & \(\valNTest\) rótulos & Métricas finais \\
\bottomrule
\end{tabular}
\end{table}

Os rótulos indicam se há pulso (\texttt{event\_present}), a amostra
aproximada de início (\texttt{t0\_sample}, igual a \(-1\) quando não há
evento) e a amplitude verdadeira (\texttt{amplitude\_A}).

\subsection{Procedimento}

A análise seguiu a ordem abaixo, com verificação em cada etapa antes de avançar:

\begin{enumerate}
  \item caracterização do ruído (média, variância, autocorrelação e PSD);
  \item análise temporal e espectral do template;
  \item projeto e aplicação do filtro de Wiener;
  \item filtro casado e combinação Wiener\,+\,casado;
  \item estimação de amplitude por BLUE (branco e colorido);
  \item limiar com FA aproximada de \(1\%\) nas janelas de ruído puro;
  \item avaliação completa no treino e escolha dos métodos;
  \item avaliação única no teste, com parâmetros congelados.
\end{enumerate}

Os métodos mínimos exigidos pelo enunciado foram implementados: sinal bruto,
Wiener, filtro casado, BLUE e a combinação Wiener\,+\,casado. Cada linha das
tabelas de resultados corresponde a um \emph{pipeline} completo:

\begin{itemize}
  \item \textbf{Sinal bruto} e \textbf{Wiener}: detecção pelo máximo de
  \(|y[n]|\) na janela (bruta ou filtrada), posição pelo índice do máximo
  (corrigido pelo pico do template) e amplitude pelo próprio valor máximo ---
  um baseline ingênuo, sem uso do template. Como estimativa da forma de onda,
  usa-se a própria janela (bruta ou filtrada).
  \item \textbf{Casado} (isolado ou após Wiener): detecção pelo máximo da
  correlação com o template, posição \(\hat{n}_0\) pelo argumento desse máximo
  e amplitude pelo estimador BLUE (branco ou colorido) aplicado à observação
  original. A forma de onda estimada é \(\hat{A}\,s[n-\hat{n}_0]\).
\end{itemize}

Nas tabelas, a coluna ``FA emp.'' é a taxa de falso alarme empírica no próprio
conjunto avaliado, isto é, a fração de janelas sem pulso classificadas como
contendo pulso pelo limiar fixado no treino.
